\scanpage{147}
\mspage{5}
lui-même était une hot-équiv. stricte, au lieu de $\alpha^\circ$, devrions avoir par $\mathrm{Dér}_5$ que pour $F\in\unsure{\mathbb D|D}(X)$, $F\xrightarrow{\sim}\alpha_*\alpha^*F$ ($\mathbb D$ dériv. à g.), ce qui est bien plus fort que $\Gamma_{\mathbb D}(X,F)\xrightarrow[\sim]{\alpha^*}\Gamma_{\mathbb D}(X',\alpha^*F)$. \deleted{Mais} La condition duale implique (si $\mathbb D$ dér. à dr.) $\alpha_!\alpha^*F\simeq F$. Mais est-il vraisemblable de s'attendre qu'elle implique aussi la condition ``\deleted{de} \unsure{du côté|nature} cohomol\textsuperscript{og.}, au lieu de l'homologique'',
\[
 \Gamma_{\mathbb D}(X,F)\xrightarrow{\sim}\Gamma_{\mathbb D}(\unsure{x'|X'},\alpha^*(F))\quad ?
\]
Le premier cas, qui devrait emporter conviction, est celui \unsure{des|du} \deleted{cat. dérivé} \inserted{dérivateur} $\mathrm{Der\,Ab}$, issu des catégories dérivées de faisceaux abéliens :

\noindent\begin{tabular}{||p{.89\linewidth}}
\uline{Question} Soit $\alpha:\unsure{X'|X}\unsure{-|\to}X$ une hot-équivalence co-stricte, $F$ un préfaisceau abélien sur $X$. A-t-on alors $H^i(X,F)\xrightarrow[\sim]{\alpha^*}H^i(X',\alpha^*\unsure{(F)|,F,|F_1})$ pour tout $i\in\mathbb Z$ ?
\end{tabular}

Si $X$ admet un objet \deleted{final} initial \inserted{$\varnothing$}, alors $\,{}_{\varnothing\backslash}X'\simeq X'$, donc $X'$ contractile, donc c'est OK pour $F$ constant.

\uline{Remarque} \unsure{Je m'aperçois} que dans la prop 1, au lieu de la condition \inserted{ou} que les $\,{}_{x\backslash}C'$ soient \inserted{fbt} filtrantes à g, il suffirait qu'elles soient non vides et \deleted{\unsure{leur $\pi_0$}} connexes, $0$-connexes.\textsuperscript{*}
% slanted
\margin{Cela signifie aussi que $\alpha$ est $W_0$-coasphérique. Pour un \unsure{traitement} cf prop. 3.2 page 43}
Le fait que les $\,{}_{x\backslash}C'$ non vides implique que $\varprojlim_C\longrightarrow\varprojlim_{C'}$ soit \uline{injectif}. Pour la surjectivité, si $\zeta'\in\varprojlim_{C'}$ (cas ensembliste), on définit $\zeta\in\varprojlim$ par les $\zeta(x)\in F(x)$ à l'aide de
% archive sheet number: 147
